\chapter{The Fail-Closed Boundary}

The parser refuses to emit a row it does not understand. The cost
of a silently wrong row is unbounded (it propagates into
aggregates, treemaps, and user actions); the cost of hard-stopping
with a typed reason is bounded (the caller dispatches to the
fallback). The discipline is concentrated in two chokepoints.

\section{The Two Chokepoints}

\begin{description}
  \item[Object headers.] One function
    (\texttt{validate\_object\_block} in the Rust crate) validates
    every block before any reader trusts its body. It checks the
    Fletcher-64 checksum, the object type, the expected storage
    class, the encryption flag, the no-header flag, optional
    \texttt{o\_oid == paddr} for physical objects, and the optional
    max-XID guard.
  \item[Record bodies.] One function (\texttt{decode\_fs\_record})
    validates every FS-tree leaf record against the rules below
    before any emission path trusts the decoded values.
\end{description}

No second path bypasses either chokepoint: a block that fails
header validation never reaches a body decoder; a record that
fails body validation never reaches the namespace builder. A
representative hard stop, from the body decoder:

\begin{minted}{rust}
if blob_header.xf_used_data as usize != padded_total {
    return Err(ScanError::InvalidObject(format!(
        "xf_used_data {} disagrees with padded value-area sum {}",
        blob_header.xf_used_data, padded_total,
    )));
}
\end{minted}

The error carries the rule that failed, the observed value, and
the expected value. The caller dispatches on the variant
(\texttt{InvalidObject}); the substring is preserved for
diagnostic display.

\section{The Per-Record Catalogue}

\textbf{Observation.} Twenty-one per-record fail-closed cases are
enforced by Rust unit tests on synthetic raw blocks. Each test
asserts that the body decoder returns a typed
\texttt{InvalidObject(reason)} with the rule's substring.

\subsection{Short bodies}

\begin{enumerate}
  \item Inode value shorter than the fixed \texttt{j\_inode\_val\_t}
    struct.
  \item Directory-record value shorter than the fixed
    \texttt{j\_drec\_val\_t} struct.
  \item Sibling-link value with \texttt{name\_len} exceeding the bytes
    actually available.
  \item Sibling-map value shorter than 8 bytes.
\end{enumerate}

\subsection{Malformed names}

\begin{enumerate}
  \setcounter{enumi}{4}
  \item Variable-length name length longer than the bytes available.
  \item Embedded NUL inside a required name (other than the
    null-terminator).
  \item Non-UTF-8 bytes in a required name.
  \item Directory-entry type outside the POSIX \texttt{DT\_*}
    allowlist.
\end{enumerate}

\subsection{Malformed xfields}

\begin{enumerate}
  \setcounter{enumi}{8}
  \item Duplicate xfield types inside one blob.
  \item \texttt{xf\_used\_data} disagreeing with
    \texttt{sum(round\_up(x\_size, 8))}.
  \item Xfield value cursor running past the blob.
  \item Xfield metadata table running past the blob.
  \item Xfield blob shorter than the 4-byte \texttt{xf\_blob\_t}
    header.
  \item Required xfield with the wrong \texttt{x\_size}
    (\texttt{DSTREAM} of length other than 40).
  \item Required xfield with the wrong \texttt{x\_size}
    (\texttt{SIBLING\_ID} of length other than 8).
\end{enumerate}

\subsection{Malformed xattrs}

\begin{enumerate}
  \setcounter{enumi}{15}
  \item Xattr value shorter than the fixed header.
  \item Xattr sets both \texttt{EMBEDDED} and \texttt{STREAM} flags.
  \item Xattr sets neither flag.
  \item Xattr sets unknown flag bits.
  \item Xattr body length disagrees with \texttt{xdata\_len}.
  \item Xattr stream body shorter than \texttt{j\_xattr\_dstream\_t}.
\end{enumerate}

\section{The Cross-Record Cases}

Three SR-016 cases depend on consistency across more than one record
and therefore cannot be enforced by synthetic single-record tests.
They are deferred to fixture-level probes; the open questions chapter
tracks their status.

\begin{description}
  \item[Drec type vs inode mode.] A directory entry's
    \texttt{DT\_*} type and the referenced inode's \texttt{S\_IFMT}
    mode must agree. A regular-file inode pointed at by a
    \texttt{DT\_DIR} entry is malformed at the volume level even
    though both records validate in isolation.
  \item[Sibling-map closure.] A directory entry that carries
    \texttt{DREC\_EXT\_TYPE\_SIBLING\_ID} requires a SIBLING\_MAP
    record that maps that sibling ID to the entry's child file ID.
    A missing closure produces a hard-link group with a dangling
    edge.
  \item[Compressed-size precedence conflict.] An inode that sets
    \texttt{INODE\_HAS\_UNCOMPRESSED\_SIZE} but does not carry a
    \texttt{com.apple.decmpfs} xattr is internally inconsistent at
    the cross-record level. Chapter~8 develops the precedence rule
    that prevents the conflict from producing wrong rows.
\end{description}

These cases will close as the project lands EX-22 and the future
cross-record consistency probe.

\section{Why The Catalogue Is Worth Its Weight}

The catalogue exercises in three not-ordinary cases:

\begin{description}
  \item[Format drift.] APFS reserves bits and types over time. A
    container written by a newer macOS may carry an xattr flag, an
    xfield type, or an incompatible-feature bit the parser does not
    recognise. The unknown-bit hard stops mean a newer container
    produces a typed refusal rather than a quietly wrong row.
  \item[Genuinely malformed images.] Recovery scenarios, partial
    image dumps, or test corpora may include containers with
    structurally inconsistent records. The catalogue refuses them
    with a precise reason, which the caller can present as a
    diagnostic.
  \item[Parser bugs.] If a future change to the body decoder
    silently re-introduces an off-by-one in the xfield cursor, the
    \texttt{xf\_used\_data} structural assertion catches it on the
    first record with a non-multiple-of-eight \texttt{x\_size}. The
    catalogue is also a regression suite.
\end{description}

\section{Object-Header Storage Classes}

\texttt{validate\_object\_block} takes the expected storage as a
parameter; a block validated under the wrong class is rejected.

\begin{description}
  \item[\texttt{Physical}.] The object's \texttt{o\_oid} equals its
    physical block address. Container superblocks, checkpoint maps,
    OMAP nodes, and free-queue nodes are physical.
  \item[\texttt{Virtual}.] The object lives in a virtual address
    space resolved by an OMAP. \texttt{o\_oid} is independent of the
    block's paddr; the validator checks
    \texttt{o\_oid == omap\_lookup.oid} instead. Volume superblocks
    (looked up through the container OMAP), FS-tree roots and
    internal nodes (looked up through the volume OMAP), and dstream
    extent nodes are virtual.
  \item[\texttt{Ephemeral}.] In-memory between transactions; reached
    through the checkpoint map. The parser uses the checkpoint map's
    mapping, never an ephemeral storage flag.
\end{description}

Mixing classes is a structural mistake. Chapter~3 records that
checkpoint-map blocks are themselves physical despite naming
ephemeral objects; chapter~5 records that FS-tree roots are virtual
despite living at a known paddr. Both rules are isolated to their
call sites.

When the boundary refuses a source, the caller dispatches to the
POSIX-traversal fallback (chapter~11), which produces the same row
shape through supported APIs.
